\documentclass{amsart}
\usepackage[UTF8,fontset=windows]{ctex}
\usepackage[a4paper,margin=24mm]{geometry}
% For dealing with references we use the comment environment
\usepackage{verbatim}
\newenvironment{reference}{\comment}{\endcomment}
%\newenvironment{reference}{}{}
\newenvironment{slogan}{\comment}{\endcomment}
\newenvironment{history}{\comment}{\endcomment}

% For commutative diagrams we use Xy-pic
\usepackage[all]{xy}

% We use 2cell for 2-commutative diagrams.
\xyoption{2cell}
\UseAllTwocells

% We use multicol for the list of chapters between chapters
\usepackage{multicol}

% This is generally recommended for better output
\usepackage{lmodern}
\usepackage[T1]{fontenc}

% For cross-file-references

% Package for hypertext links:
\usepackage{hyperref}

% For any local file, say "hello.tex" you want to link to please

% Theorem environments.
%
\theoremstyle{plain}
\newtheorem{theorem}[subsection]{Theorem}
\newtheorem{proposition}[subsection]{Proposition}
\newtheorem{lemma}[subsection]{Lemma}

\theoremstyle{definition}
\newtheorem{definition}[subsection]{Definition}
\newtheorem{example}[subsection]{Example}
\newtheorem{exercise}[subsection]{Exercise}
\newtheorem{situation}[subsection]{Situation}

\theoremstyle{remark}
\newtheorem{remark}[subsection]{Remark}
\newtheorem{remarks}[subsection]{Remarks}

\numberwithin{equation}{subsection}

% Macros
%
\def\lim{\mathop{\mathrm{lim}}\nolimits}
\def\colim{\mathop{\mathrm{colim}}\nolimits}
\def\Spec{\mathop{\mathrm{Spec}}}
\def\Hom{\mathop{\mathrm{Hom}}\nolimits}
\def\Ext{\mathop{\mathrm{Ext}}\nolimits}
\def\SheafHom{\mathop{\mathcal{H}\!\mathit{om}}\nolimits}
\def\SheafExt{\mathop{\mathcal{E}\!\mathit{xt}}\nolimits}
\def\Sch{\mathit{Sch}}
\def\Mor{\mathop{\mathrm{Mor}}\nolimits}
\def\Ob{\mathop{\mathrm{Ob}}\nolimits}
\def\Sh{\mathop{\mathit{Sh}}\nolimits}
\def\NL{\mathop{N\!L}\nolimits}
\def\CH{\mathop{\mathrm{CH}}\nolimits}
\def\proetale{{pro\text{-}\acute{e}tale}}
\def\etale{{\acute{e}tale}}
\def\QCoh{\mathit{QCoh}}
\def\Ker{\mathop{\mathrm{Ker}}}
\def\Im{\mathop{\mathrm{Im}}}
\def\Coker{\mathop{\mathrm{Coker}}}
\def\Coim{\mathop{\mathrm{Coim}}}

% Boxtimes
%
\DeclareMathSymbol{\boxtimes}{\mathbin}{AMSa}{"02}

%
% Macros for moduli stacks/spaces
%
\def\QCohstack{\mathcal{QC}\!\mathit{oh}}
\def\Cohstack{\mathcal{C}\!\mathit{oh}}
\def\Spacesstack{\mathcal{S}\!\mathit{paces}}
\def\Quotfunctor{\mathrm{Quot}}
\def\Hilbfunctor{\mathrm{Hilb}}
\def\Curvesstack{\mathcal{C}\!\mathit{urves}}
\def\Polarizedstack{\mathcal{P}\!\mathit{olarized}}
\def\Complexesstack{\mathcal{C}\!\mathit{omplexes}}
% \Pic is the operator that assigns to X its picard group, usage \Pic(X)
% \Picardstack_{X/B} denotes the Picard stack of X over B
% \Picardfunctor_{X/B} denotes the Picard functor of X over B
\def\Pic{\mathop{\mathrm{Pic}}\nolimits}
\def\Picardstack{\mathcal{P}\!\mathit{ic}}
\def\Picardfunctor{\mathrm{Pic}}
\def\Deformationcategory{\mathcal{D}\!\mathit{ef}}

\hypersetup{hidelinks}
\widowpenalty=10000\clubpenalty=10000\displaywidowpenalty=10000
\begin{document}
\title[Stacks Project, Tag 0DIW]{平方零厚化上导出复形的障碍消失判据}
\maketitle
\noindent\textbf{Counting scope.} The primary problem is Tag 0DIW, marked below. All preceding numbered statements are supporting author context and are not counted as additional problems. This is a local pilot with a temporary ID, not production admission.


\noindent Copyright (C) 2005--2025 Johan de Jong. Stacks Project authors. Source: \href{https://stacks.math.columbia.edu/tag/0DIW}{Tag 0DIW}.

\noindent Difficulty H2: 本题证明导出复形提升存在当且仅当其规范二次障碍类为零。作者把零同伦作为校正项在纤维积模上定义新的微分，再通过导出张量、上同调下降归纳和四引理证明模掉平方零理想后得到原复形。规范类的构造可以明确作为背景，但校正复形与拟同构验证不能由背景拿掉，判H2。.

\noindent History: excerpt from revision local cached snapshot; exact chapter hashes; revision unknown, defos.tex, lines 5357--5495. Only reference presentation was adapted. Licensed under GNU FDL 1.2 or later; no invariant sections or cover texts. See accompanying COPYING.

\begin{lemma}
\label{lemma-canonical-class}
Let $\mathcal{C}$ be a site. Let $\mathcal{O} \to \mathcal{O}_0$
be a surjection of sheaves of rings whose kernel is an ideal sheaf
$\mathcal{I}$ of square zero. For every object
$K_0$ in $D^-(\mathcal{O}_0)$ there is a canonical map
$$
\omega(K_0) :
K_0 \longrightarrow
K_0 \otimes_{\mathcal{O}_0}^\mathbf{L} \mathcal{I}[2]
$$
in $D(\mathcal{O}_0)$ such that for any map
$K_0 \to L_0$ in $D^-(\mathcal{O}_0)$ the diagram
$$
\xymatrix{
K_0 \ar[d] \ar[rr]_-{\omega(K_0)} & &
(K_0 \otimes^\mathbf{L}_{\mathcal{O}_0} \mathcal{I})[2] \ar[d] \\
L_0 \ar[rr]^-{\omega(L_0)} & &
(L_0 \otimes^\mathbf{L}_{\mathcal{O}_0} \mathcal{I})[2]
}
$$
commutes.
\end{lemma}

\begin{proof}
Represent $K_0$ by any complex
$\mathcal{K}_0^\bullet$ of $\mathcal{O}_0$-modules.
Apply Lemma \ref{lemma-lift-complex}
with $\mathcal{G}^n = 0$ for all $n$.
Denote $d : \mathcal{F}^n \to \mathcal{F}^{n + 1}$
the maps produced by the lemma. Then we see that
$d \circ d : \mathcal{F}^n \to \mathcal{F}^{n + 2}$
is zero modulo $\mathcal{I}$. Since $\mathcal{F}^n$ is flat,
we see that
$\mathcal{I}\mathcal{F}^n =
\mathcal{F}^n \otimes_{\mathcal{O}} \mathcal{I} =
\mathcal{F}^n_0 \otimes_{\mathcal{O}_0} \mathcal{I}$.
Hence we obtain a canonical map of complexes
$$
d \circ d : \mathcal{F}_0^\bullet
\longrightarrow
(\mathcal{F}_0^\bullet \otimes_{\mathcal{O}_0} \mathcal{I})[2]
$$
Since $\mathcal{F}_0^\bullet$ is a bounded above complex
of flat $\mathcal{O}_0$-modules, it is K-flat and may be used
to compute derived tensor product. Moreover, the map of
complexes $\mathcal{F}_0^\bullet \to \mathcal{K}_0^\bullet$
is a quasi-isomorphism by construction. Therefore the source and
target of the map just constructed represent $K_0$ and
$K_0 \otimes_{\mathcal{O}_0}^\mathbf{L} \mathcal{I}[2]$
and we obtain our map $\omega(K_0)$.

\medskip\noindent
Let us show that this procedure is compatible with maps of complexes.
Namely, let $\mathcal{L}_0^\bullet$ represent another object of
$D^-(\mathcal{O}_0)$ and suppose that
$$
\mathcal{K}_0^\bullet \longrightarrow \mathcal{L}_0^\bullet
$$
is a map of complexes. Apply Lemma \ref{lemma-lift-complex}
for the complex $\mathcal{L}_0^\bullet$, the flat modules
$\mathcal{F}^n$, the maps $\mathcal{F}^n \to \mathcal{F}^{n + 1}$,
and the compositions
$\mathcal{F}^n \to \mathcal{K}_0^n \to \mathcal{L}_0^n$
(we apologize for the reversal of letters used).
We obtain flat modules $\mathcal{G}^n$, maps
$\mathcal{F}^n \to \mathcal{G}^n$, maps
$\mathcal{G}^n \to \mathcal{G}^{n + 1}$, and maps
$\mathcal{G}^n \to \mathcal{L}_0^n$ with all properties
as in the lemma. Then it is clear that
$$
\xymatrix{
\mathcal{F}_0^\bullet \ar[d] \ar[r] &
(\mathcal{F}_0^\bullet \otimes_{\mathcal{O}_0} \mathcal{I})[2] \ar[d] \\
\mathcal{G}_0^\bullet \ar[r] &
(\mathcal{G}_0^\bullet \otimes_{\mathcal{O}_0} \mathcal{I})[2]
}
$$
is a commutative diagram of complexes.

\medskip\noindent
To see that $\omega(K_0)$ is well defined, suppose that we have two complexes
$\mathcal{K}_0^\bullet$ and $(\mathcal{K}'_0)^\bullet$
of $\mathcal{O}_0$-modules representing $K_0$ and two systems
$(\mathcal{F}^n, d : \mathcal{F}^n \to \mathcal{F}^{n + 1},
\mathcal{F}^n \to \mathcal{K}_0^n)$
and
$((\mathcal{F}')^n, d : (\mathcal{F}')^n \to (\mathcal{F}')^{n + 1},
(\mathcal{F}')^n \to \mathcal{K}_0^n)$
as above. Then we can choose a complex $(\mathcal{K}''_0)^\bullet$
and quasi-isomorphisms
$\mathcal{K}_0^\bullet  \to (\mathcal{K}''_0)^\bullet$
and
$(\mathcal{K}'_0)^\bullet  \to (\mathcal{K}''_0)^\bullet$
realizing the fact that both complexes represent $K_0$ in the
derived category. Next, we apply the result of the previous paragraph
to
$$
(\mathcal{K}_0)^\bullet \oplus (\mathcal{K}'_0)^\bullet
\longrightarrow
(\mathcal{K}''_0)^\bullet
$$
This produces a commutative diagram
$$
\xymatrix{
\mathcal{F}_0^\bullet \oplus (\mathcal{F}'_0)^\bullet
\ar[d] \ar[r] &
(\mathcal{F}_0^\bullet \otimes_{\mathcal{O}_0} \mathcal{I})[2] \oplus
((\mathcal{F}'_0)^\bullet \otimes_{\mathcal{O}_0} \mathcal{I})[2] \ar[d] \\
\mathcal{G}_0^\bullet \ar[r] &
(\mathcal{G}_0^\bullet \otimes_{\mathcal{O}_0} \mathcal{I})[2]
}
$$
Since the vertical arrows give quasi-isomorphisms on the summands
we conclude the desired commutativity in $D(\mathcal{O}_0)$.

\medskip\noindent
Having established well-definedness, the statement on compatibility
with maps is a consequence of the result in the second
paragraph.
\end{proof}


\begin{lemma}
\label{lemma-lift-complex}
Let $\mathcal{C}$ be a site. Let $\mathcal{O} \to \mathcal{O}_0$
be a surjection of sheaves of rings. Assume given the following data
\begin{enumerate}
\item flat $\mathcal{O}$-modules $\mathcal{G}^n$,
\item maps of $\mathcal{O}$-modules $\mathcal{G}^n \to \mathcal{G}^{n + 1}$,
\item a complex $\mathcal{K}_0^\bullet$ of $\mathcal{O}_0$-modules,
\item maps of $\mathcal{O}$-modules $\mathcal{G}^n \to \mathcal{K}_0^n$
\end{enumerate}
such that
\begin{enumerate}
\item[(a)] $H^n(\mathcal{K}_0^\bullet) = 0$ for $n \gg 0$,
\item[(b)] $\mathcal{G}^n = 0$ for $n \gg 0$,
\item[(c)] with
$\mathcal{G}^n_0 = \mathcal{G}^n \otimes_\mathcal{O} \mathcal{O}_0$
the induced maps determine a complex $\mathcal{G}_0^\bullet$ and a map
of complexes $\mathcal{G}_0^\bullet \to \mathcal{K}_0^\bullet$.
\end{enumerate}
Then there exist
\begin{enumerate}
\item[(\romannumeral1)]
flat $\mathcal{O}$-modules $\mathcal{F}^n$,
\item[(\romannumeral2)]
maps of $\mathcal{O}$-modules $\mathcal{F}^n \to \mathcal{F}^{n + 1}$,
\item[(\romannumeral3)]
maps of $\mathcal{O}$-modules $\mathcal{F}^n \to \mathcal{K}_0^n$,
\item[(\romannumeral4)]
maps of $\mathcal{O}$-modules $\mathcal{G}^n \to \mathcal{F}^n$,
\end{enumerate}
such that $\mathcal{F}^n = 0$ for $n \gg 0$, such that the diagrams
$$
\xymatrix{
\mathcal{G}^n \ar[r] \ar[d] & \mathcal{G}^{n + 1} \ar[d] \\
\mathcal{F}^n \ar[r] & \mathcal{F}^{n + 1}
}
$$
commute for all $n$, such that the composition
$\mathcal{G}^n \to \mathcal{F}^n \to \mathcal{K}_0^n$
is the given map $\mathcal{G}^n \to \mathcal{K}_0^n$, and such that with
$\mathcal{F}^n_0 = \mathcal{F}^n \otimes_\mathcal{O} \mathcal{O}_0$
we obtain a complex $\mathcal{F}_0^\bullet$ and map of complexes
$\mathcal{F}_0^\bullet \to \mathcal{K}_0^\bullet$ which is a
quasi-isomorphism.
\end{lemma}

\begin{proof}
We will prove by descending induction on $e$ that we can find $\mathcal{F}^n$,
$\mathcal{G}^n \to \mathcal{F}^n$, and
$\mathcal{F}^n \to \mathcal{F}^{n + 1}$ for $n \geq e$
fitting into a commutative diagram
$$
\xymatrix{
\ldots \ar[r] &
\mathcal{G}^{e - 1} \ar[r] \ar@/_2pc/[dd] &
\mathcal{G}^e \ar[d] \ar[r] \ar@/_2pc/[dd] &
\mathcal{G}^{e + 1} \ar[d] \ar[r] \ar@/_2pc/[dd]|\hole &
\ldots \\
& &
\mathcal{F}^e \ar[d] \ar[r] &
\mathcal{F}^{e + 1} \ar[d] \ar[r] & \ldots \\
\ldots  \ar[r] &
\mathcal{K}_0^{e - 1} \ar[r] &
\mathcal{K}_0^e \ar[r] &
\mathcal{K}_0^{e + 1} \ar[r] & \ldots
}
$$
such that $\mathcal{F}_0^\bullet$ is a complex,
the induced map $\mathcal{F}_0^\bullet \to \mathcal{K}_0^\bullet$
induces an isomorphism on $H^n$ for $n > e$ and a surjection
for $n = e$. For $e \gg 0$ this is true because we can take
$\mathcal{F}^n = 0$ for $n \geq e$ in that case by assumptions
(a) and (b).

\medskip\noindent
Induction step. We have to construct $\mathcal{F}^{e - 1}$
and the maps $\mathcal{G}^{e - 1} \to \mathcal{F}^{e - 1}$,
$\mathcal{F}^{e - 1} \to \mathcal{F}^e$, and
$\mathcal{F}^{e - 1} \to \mathcal{K}_0^{e - 1}$.
We will choose $\mathcal{F}^{e - 1} = A \oplus B \oplus C$
as a direct sum of three pieces.

\medskip\noindent
For the first we take $A = \mathcal{G}^{e - 1}$ and we choose our map
$\mathcal{G}^{e - 1} \to \mathcal{F}^{e - 1}$ to be the inclusion of
the first summand. The maps $A \to \mathcal{K}^{e - 1}_0$
and $A \to \mathcal{F}^e$ will be the obvious ones.

\medskip\noindent
To choose $B$ we consider the surjection (by induction hypothesis)
$$
\gamma :
\Ker(\mathcal{F}^e_0 \to \mathcal{F}^{e + 1}_0)
\longrightarrow
\Ker(\mathcal{K}^e_0 \to \mathcal{K}^{e + 1}_0)/
\Im(\mathcal{K}^{e - 1}_0 \to \mathcal{K}^e_0)
$$
We can choose a set $I$, for each $i \in I$
an object $U_i$ of $\mathcal{C}$, and sections
$s_i \in \mathcal{F}^e(U_i)$, $t_i \in \mathcal{K}^{e - 1}_0(U_i)$
such that
\begin{enumerate}
\item $s_i$ maps to a section of $\Ker(\gamma) \subset
\Ker(\mathcal{F}^e_0 \to \mathcal{F}^{e + 1}_0)$,
\item $s_i$ and $t_i$ map to the same section of
$\mathcal{K}^e_0$,
\item the sections $s_i$ generate $\Ker(\gamma)$ as an $\mathcal{O}_0$-module.
\end{enumerate}
We omit giving the full justification for this;
one uses that $\mathcal{F}^e \to \mathcal{F}^e_0$
is a surjective maps of sheaves of sets. Then we set
to put
$$
B = \bigoplus\nolimits_{i \in I} j_{U_i!}\mathcal{O}_{U_i}
$$
and define the maps $B \to \mathcal{F}^e$ and $B \to \mathcal{K}_0^{e - 1}$
by using $s_i$ and $t_i$ to determine where to send the summand
$j_{U_i!}\mathcal{O}_{U_i}$.

\medskip\noindent
With $\mathcal{F}^{e - 1} = A \oplus B$ and maps as above,
this produces a diagram as above for $e - 1$ such that
$\mathcal{F}_0^\bullet \to \mathcal{K}_0^\bullet$
induces an isomorphism on $H^n$ for $n \geq e$.
To get the map to be surjective on $H^{e - 1}$ we choose
the summand $C$ as follows.
Choose a set $J$, for each $j \in J$ an object $U_j$ of $\mathcal{C}$
and a section $t_j$ of $\Ker(\mathcal{K}^{e - 1}_0 \to \mathcal{K}^e_0)$
over $U_j$ such that these sections generate this kernel over
$\mathcal{O}_0$. Then we put
$$
C = \bigoplus\nolimits_{j \in J} j_{U_j!}\mathcal{O}_{U_j}
$$
and the zero map $C \to \mathcal{F}^e$ and the map
$C \to \mathcal{K}_0^{e - 1}$ by using $s_j$ to determine where to the summand
$j_{U_j!}\mathcal{O}_{U_j}$. This finishes the induction step
by taking $\mathcal{F}^{e - 1} = A \oplus B \oplus C$ and
maps as indicated.
\end{proof}


\begin{lemma}
\label{lemma-induced-map}
Let $(\mathcal{C}, \mathcal{O})$ be a ringed site.
Let $\alpha : K \to L$ be a map of $D^-(\mathcal{O})$.
Let $\mathcal{F}$ be a sheaf of $\mathcal{O}$-modules.
Let $n \in \mathbf{Z}$.
\begin{enumerate}
\item If $H^i(\alpha)$ is an isomorphism for $i \geq n$,
then $H^i(\alpha \otimes_\mathcal{O}^\mathbf{L} \text{id}_\mathcal{F})$
is an isomorphism for $i \geq n$.
\item If $H^i(\alpha)$ is an isomorphism for $i > n$ 
and surjective for $i = n$,
then $H^i(\alpha \otimes_\mathcal{O}^\mathbf{L} \text{id}_\mathcal{F})$
is an isomorphism for $i > n$ and surjective for $i = n$.
\end{enumerate}
\end{lemma}

\section*{Primary problem: Tag 0DIW}
\begin{lemma}
\label{lemma-canonical-class-obstruction}
Let $\mathcal{C}$ be a site. Let $\mathcal{O} \to \mathcal{O}_0$
be a surjection of sheaves of rings whose kernel is an ideal sheaf
$\mathcal{I}$ of square zero. For every object
$K_0$ in $D^-(\mathcal{O}_0)$ the following are equivalent
\begin{enumerate}
\item  the class
$\omega(K_0) \in
\Ext^2_{\mathcal{O}_0}(K_0, K_0 \otimes_{\mathcal{O}_0} \mathcal{I})$
constructed in Lemma \ref{lemma-canonical-class} is zero,
\item there exists $K \in D^-(\mathcal{O})$ with
$K \otimes_\mathcal{O}^\mathbf{L} \mathcal{O}_0 = K_0$
in $D(\mathcal{O}_0)$.
\end{enumerate}
\end{lemma}

\begin{proof}
Let $K$ be as in (2). Then we can represent $K$ by a bounded above
complex $\mathcal{F}^\bullet$ of flat $\mathcal{O}$-modules.
Then $\mathcal{F}_0^\bullet =
\mathcal{F}^\bullet \otimes_{\mathcal{O}} \mathcal{O}_0$
represents $K_0$ in $D(\mathcal{O}_0)$.
Since $d_{\mathcal{F}^\bullet} \circ d_{\mathcal{F}^\bullet} = 0$
as $\mathcal{F}^\bullet$ is a complex, we see from the very construction
of $\omega(K_0)$ that it is zero.

\medskip\noindent
Assume (1). Let $\mathcal{F}^n$, $d : \mathcal{F}^n \to \mathcal{F}^{n + 1}$
be as in the construction of $\omega(K_0)$. The nullity of
$\omega(K_0)$ implies that the map
$$
\omega = d \circ d : \mathcal{F}_0^\bullet
\longrightarrow
(\mathcal{F}_0^\bullet \otimes_{\mathcal{O}_0} \mathcal{I})[2]
$$
is zero in $D(\mathcal{O}_0)$. By definition of the derived category
as the localization of the homotopy category of complexes
of $\mathcal{O}_0$-modules, there exists a quasi-isomorphism
$\alpha : \mathcal{G}_0^\bullet \to \mathcal{F}_0^\bullet$
such that there exist $\mathcal{O}_0$-modules maps
$h^n : \mathcal{G}_0^n \to
\mathcal{F}_0^{n + 1} \otimes_\mathcal{O} \mathcal{I}$
with
$$
\omega \circ \alpha =
d_{\mathcal{F}_0^\bullet \otimes \mathcal{I}} \circ h +
h \circ d_{\mathcal{G}_0^\bullet}
$$
We set
$$
\mathcal{H}^n = \mathcal{F}^n \times_{\mathcal{F}^n_0} \mathcal{G}_0^n
$$
and we define
$$
d' : \mathcal{H}^n \longrightarrow \mathcal{H}^{n + 1},\quad
(f^n, g_0^n) \longmapsto (d(f^n) - h^n(g_0^n), d(g_0^n))
$$
with obvious notation using that
$\mathcal{F}_0^{n + 1} \otimes_{\mathcal{O}_0} \mathcal{I} =
\mathcal{F}^{n + 1} \otimes_\mathcal{O} \mathcal{I} =
\mathcal{I}\mathcal{F}^{n + 1} \subset \mathcal{F}^{n + 1}$.
Then one checks $d' \circ d' = 0$ by our choice of $h^n$
and definition of $\omega$.
Hence $\mathcal{H}^\bullet$ defines an object in $D(\mathcal{O})$.
On the other hand, there is a short exact sequence of complexes
of $\mathcal{O}$-modules
$$
0 \to \mathcal{F}_0^\bullet \otimes_{\mathcal{O}_0} \mathcal{I} \to
\mathcal{H}^\bullet \to \mathcal{G}_0^\bullet \to 0
$$
We still have to show that
$\mathcal{H}^\bullet \otimes_\mathcal{O}^\mathbf{L} \mathcal{O}_0$
is isomorphic to $K_0$.
Choose a quasi-isomorphism
$\mathcal{E}^\bullet \to \mathcal{H}^\bullet$
where $\mathcal{E}^\bullet$ is a bounded above complex of flat
$\mathcal{O}$-modules. We obtain a commutative diagram
$$
\xymatrix{
0 \ar[r] &
\mathcal{E}^\bullet \otimes_\mathcal{O} \mathcal{I} \ar[d]^\beta \ar[r] &
\mathcal{E}^\bullet \ar[d]^\gamma \ar[r] &
\mathcal{E}_0^\bullet \ar[d]^\delta \ar[r] &
0 \\
0 \ar[r] &
\mathcal{F}_0^\bullet \otimes_{\mathcal{O}_0} \mathcal{I} \ar[r] &
\mathcal{H}^\bullet \ar[r] &
\mathcal{G}_0^\bullet \ar[r] &
0
}
$$
We claim that $\delta$ is a quasi-isomorphism. Since $H^i(\delta)$
is an isomorphism for $i \gg 0$, we can use descending induction
on $n$ such that $H^i(\delta)$ is an isomorphism for $i \geq n$.
Observe that
$\mathcal{E}^\bullet \otimes_\mathcal{O} \mathcal{I}$
represents
$\mathcal{E}_0^\bullet \otimes_{\mathcal{O}_0}^\mathbf{L} \mathcal{I}$,
that
$\mathcal{F}_0^\bullet \otimes_{\mathcal{O}_0} \mathcal{I}$
represents
$\mathcal{G}_0^\bullet \otimes_{\mathcal{O}_0}^\mathbf{L} \mathcal{I}$,
and that
$\beta = \delta \otimes_{\mathcal{O}_0}^\mathbf{L} \text{id}_\mathcal{I}$
as maps in $D(\mathcal{O}_0)$. This is true because
$\beta =
(\alpha \otimes \text{id}_\mathcal{I})
\circ
(\delta \otimes \text{id}_\mathcal{I})$.
Suppose that $H^i(\delta)$ is an isomorphism in degrees $\geq n$.
Then the same is true for $\beta$ by what we just said
and Lemma \ref{lemma-induced-map}.
Then we can look at the diagram
$$
\xymatrix{
H^{n - 1}(\mathcal{E}^\bullet \otimes_\mathcal{O} \mathcal{I})
\ar[r] \ar[d]^{H^{n - 1}(\beta)} &
H^{n - 1}(\mathcal{E}^\bullet) \ar[r] \ar[d] &
H^{n - 1}(\mathcal{E}_0^\bullet) \ar[r] \ar[d]^{H^{n - 1}(\delta)} &
H^n(\mathcal{E}^\bullet \otimes_\mathcal{O} \mathcal{I})
\ar[r] \ar[d]^{H^n(\beta)} &
H^n(\mathcal{E}^\bullet) \ar[d] \\
H^{n - 1}(\mathcal{F}_0^\bullet \otimes_\mathcal{O} \mathcal{I}) \ar[r] &
H^{n - 1}(\mathcal{H}^\bullet) \ar[r] &
H^{n - 1}(\mathcal{G}_0^\bullet) \ar[r] &
H^n(\mathcal{F}_0^\bullet \otimes_\mathcal{O} \mathcal{I}) \ar[r] &
H^n(\mathcal{H}^\bullet)
}
$$
Using Homology, Lemma \href{https://stacks.math.columbia.edu/tag/05QA}{lemma four lemma (Tag 05QA)}
we see that $H^{n - 1}(\delta)$ is surjective.
This in turn implies that $H^{n - 1}(\beta)$ is surjective
by Lemma \ref{lemma-induced-map}.
Using Homology, Lemma \href{https://stacks.math.columbia.edu/tag/05QA}{lemma four lemma (Tag 05QA)}
again we see that $H^{n - 1}(\delta)$ is an isomorphism.
The claim holds by induction, so $\delta$ is a quasi-isomorphism
which is what we wanted to show.
\end{proof}
\end{document}
